\section{Entity-Graph (5123034, 5123060, 5123097)}
\label{sec:performance_entity_graphs}

Entity Graphs in JPA/Hibernate ermöglichen es, gezielt zu steuern, welche Assoziationen beim Laden einer Entität mitgeladen werden. 
Ohne diese Kontrolle entstehen zwei klassische Probleme: das $N+1$-Problem und unnötiges Lazy-Loading.

\subsection*{Das $N+1$-Problem}
Ohne Entity Graphs lädt Hibernate verknüpfte Entitäten standardmäßig einzeln pro Datensatz (\textit{lazy}). 
Beispielsweise bei einer Liste von Ärzten würde Hibernate zunächst eine Abfrage für die Ärzte ausführen, danach aber für jeden Arzt eine separate Abfrage für dessen Zuweisungen, Diagnosen oder Medikamente.
Mit einem Entity Graph werden alle benötigten Assoziationen in einer einzigen \texttt{JOIN}-Abfrage geladen. 
Dadurch reduziert sich der Overhead durch Netzwerk-Roundtrips zur Datenbank drastisch.

\subsection*{Anwendung im Projekt}
In Tabelle~\ref{tab:entity_graphs_anwendung} sind die konkreten Problemstellen unseres Hospital Management Systems und deren jeweilige Lösung durch Entity Graphs gegenübergestellt.

\begin{table}[H]
\centering
\small
\caption{Übersicht der Performance-Optimierung durch Entity Graphs}
\label{tab:entity_graphs_anwendung}
\begin{tabularx}{\textwidth}{l X X}
\hline
\textbf{Entität} & \textbf{Problemstelle ohne Graph} & \textbf{Lösung durch Graph} \\ \hline
Doctor & Abteilung \& Krankenhausinfo separat geladen & Alles in einem \texttt{JOIN} \\
Nurse & Schichten und Zuweisungen ($N+1$) & \textit{Eager-Join} auf relevante Felder \\
Diagnosis & Zugehörige Booking \& Patient \textit{lazy} nachgeladen & Graph lädt Kontext direkt mit \\
Medication & Diagnose-Verknüpfung erst bei Zugriff geladen & Direkt verfügbar ohne Extra-Query \\ \hline
\end{tabularx}
\end{table}

\subsection*{Ergebnis}
Durch den gezielten Einsatz von \texttt{@EntityGraph} in den Repository-Methoden reduziert sich die Anzahl der Datenbankabfragen erheblich, insbesondere bei Listenabfragen über viele Datensätze. 
Gleichzeitig wird \textit{Overfetching} vermieden, da nur die tatsächlich benötigten Assoziationen mitgeladen werden, nicht pauschal alles.
Das führt zu einer messbaren Verbesserung bei den Antwortzeiten der API, einer geringeren Datenbanklast und einem stabileren Verhalten unter hoher Last.
Die Benchmarks der nachfolgenden Section enthalten bereits diese Verbesserungen.